%%
%% Automatically generated file from DocOnce source
%% (https://github.com/doconce/doconce/)
%% doconce format latex week12.do.txt --minted_latex_style=trac --latex_admon=paragraph --no_mako
%%
% #ifdef PTEX2TEX_EXPLANATION
%%
%% The file follows the ptex2tex extended LaTeX format, see
%% ptex2tex: https://code.google.com/p/ptex2tex/
%%
%% Run
%%      ptex2tex myfile
%% or
%%      doconce ptex2tex myfile
%%
%% to turn myfile.p.tex into an ordinary LaTeX file myfile.tex.
%% (The ptex2tex program: https://code.google.com/p/ptex2tex)
%% Many preprocess options can be added to ptex2tex or doconce ptex2tex
%%
%%      ptex2tex -DMINTED myfile
%%      doconce ptex2tex myfile envir=minted
%%
%% ptex2tex will typeset code environments according to a global or local
%% .ptex2tex.cfg configure file. doconce ptex2tex will typeset code
%% according to options on the command line (just type doconce ptex2tex to
%% see examples). If doconce ptex2tex has envir=minted, it enables the
%% minted style without needing -DMINTED.
% #endif

% #define PREAMBLE

% #ifdef PREAMBLE
%-------------------- begin preamble ----------------------

\documentclass[%
oneside,                 % oneside: electronic viewing, twoside: printing
final,                   % draft: marks overfull hboxes, figures with paths
10pt]{article}

\listfiles               %  print all files needed to compile this document

\usepackage{relsize,makeidx,color,setspace,amsmath,amsfonts,amssymb}
\usepackage[table]{xcolor}
\usepackage{bm,ltablex,microtype}

\usepackage[pdftex]{graphicx}

\usepackage{ptex2tex}
% #ifdef MINTED
\usepackage{minted}
\usemintedstyle{default}
% #endif

\usepackage[T1]{fontenc}
%\usepackage[latin1]{inputenc}
\usepackage{ucs}
\usepackage[utf8x]{inputenc}

\usepackage{lmodern}         % Latin Modern fonts derived from Computer Modern

% Hyperlinks in PDF:
\definecolor{linkcolor}{rgb}{0,0,0.4}
\usepackage{hyperref}
\hypersetup{
    breaklinks=true,
    colorlinks=true,
    linkcolor=linkcolor,
    urlcolor=linkcolor,
    citecolor=black,
    filecolor=black,
    %filecolor=blue,
    pdfmenubar=true,
    pdftoolbar=true,
    bookmarksdepth=3   % Uncomment (and tweak) for PDF bookmarks with more levels than the TOC
    }
%\hyperbaseurl{}   % hyperlinks are relative to this root

\setcounter{tocdepth}{2}  % levels in table of contents

% --- fancyhdr package for fancy headers ---
\usepackage{fancyhdr}
\fancyhf{} % sets both header and footer to nothing
\renewcommand{\headrulewidth}{0pt}
\fancyfoot[LE,RO]{\thepage}
% Ensure copyright on titlepage (article style) and chapter pages (book style)
\fancypagestyle{plain}{
  \fancyhf{}
  \fancyfoot[C]{{\footnotesize \copyright\ 1999-2026, Morten Hjorth-Jensen  Email morten.hjorth-jensen@fys.uio.no. Released under CC Attribution-NonCommercial 4.0 license}}
%  \renewcommand{\footrulewidth}{0mm}
  \renewcommand{\headrulewidth}{0mm}
}
% Ensure copyright on titlepages with \thispagestyle{empty}
\fancypagestyle{empty}{
  \fancyhf{}
  \fancyfoot[C]{{\footnotesize \copyright\ 1999-2026, Morten Hjorth-Jensen  Email morten.hjorth-jensen@fys.uio.no. Released under CC Attribution-NonCommercial 4.0 license}}
  \renewcommand{\footrulewidth}{0mm}
  \renewcommand{\headrulewidth}{0mm}
}

\pagestyle{fancy}


\usepackage[framemethod=TikZ]{mdframed}

% --- begin definitions of admonition environments ---

% --- end of definitions of admonition environments ---

% prevent orhpans and widows
\clubpenalty = 10000
\widowpenalty = 10000

% --- end of standard preamble for documents ---


% insert custom LaTeX commands...

\raggedbottom
\makeindex
\usepackage[totoc]{idxlayout}   % for index in the toc
\usepackage[nottoc]{tocbibind}  % for references/bibliography in the toc

%-------------------- end preamble ----------------------

\begin{document}

% matching end for #ifdef PREAMBLE
% #endif

\newcommand{\exercisesection}[1]{\subsection*{#1}}


% ------------------- main content ----------------------



% ----------------- title -------------------------

\thispagestyle{empty}

\begin{center}
{\LARGE\bf
\begin{spacing}{1.25}
Boltzmann machines and deep learning
\end{spacing}
}
\end{center}

% ----------------- author(s) -------------------------

\begin{center}
{\bf Morten Hjorth-Jensen  Email morten.hjorth-jensen@fys.uio.no}
\end{center}

    \begin{center}
% List of all institutions:
\centerline{{\small Department of Physics and Center fo Computing in Science Education, University of Oslo, Oslo, Norway}}
\end{center}
    
% ----------------- end author(s) -------------------------

% --- begin date ---
\begin{center}
April 17, 2026
\end{center}
% --- end date ---

\vspace{1cm}


% !split
\subsection{Plans for the week of April 13-17, 2026}


% --- begin paragraph admon ---
\paragraph{}
\begin{itemize}
\item Neural Networks and Boltzmann Machines, examples of generative deep learning

\item Discussions of project 2
\end{itemize}

\noindent
% --- end paragraph admon ---



% !split
\subsection{Boltzmann Machines, marginal and conditional probabilities}

A generative model can learn to represent and sample from a
probability distribution. The core idea is to learn a parametric model
of the probability distribution from which the training data was
drawn. As an example
\begin{enumerate}
\item A model for images could learn to draw new examples of cats and dogs, given a training dataset of images of cats and dogs.

\item Generate a sample of an ordered or disordered Ising model phase, having been given samples of such phases.

\item Model the trial function for Monte Carlo calculations
\end{enumerate}

\noindent
% !split
\subsection{Generative and discriminative models}

Generative and discriminative models use both gradient-descent based
learning procedures for minimizing cost functions

However, in energy based models we don't use backpropagation and automatic
differentiation for computing gradients, instead we turn to Markov
Chain Monte Carlo methods.

A typical deep neural network has several hidden layers. A restricted
Boltzmann machine has normally one hidden layer, however several RBMs can
be stacked to make up deep Belief Networks, of which they constitute
the building blocks.

% !split
\subsection{Basics of the Boltzmann machine}

% --- begin paragraph admon ---
\paragraph{}
A BM is what we would call an undirected probabilistic graphical model
with stochastic continuous or discrete units.
% --- end paragraph admon ---



% --- begin paragraph admon ---
\paragraph{}
It is interpreted as a stochastic recurrent neural network where the
state of each unit(neurons/nodes) depends on the units it is connected
to. The weights in the network represent thus the strength of the
interaction between various units/nodes.
% --- end paragraph admon ---



% !split
\subsection{More about the basics}

% --- begin paragraph admon ---
\paragraph{}
A standard BM network is divided into a set of observable and visible units $\bm{x}$ and a set of unknown hidden units/nodes $\bm{h}$.
% --- end paragraph admon ---




% --- begin paragraph admon ---
\paragraph{}
Additionally there can be bias nodes for the hidden and visible layers. These biases are normally set to $1$.
% --- end paragraph admon ---




% --- begin paragraph admon ---
\paragraph{}
BMs are stackable, meaning they cwe can train a BM which serves as input to another BM. We can construct deep networks for learning complex PDFs. The layers can be trained one after another, a feature which makes them popular in deep learning
% --- end paragraph admon ---



% !split
\subsection{Difficult to train}

However, they are often hard to train. This leads to the introduction
of so-called restricted BMs, or RBMS.  Here we take away all lateral
connections between nodes in the visible layer as well as connections
between nodes in the hidden layer.

% !split
\subsection{The network layers}

\begin{enumerate}
\item A function $\bm{x}$ that represents the visible layer, a vector of $M$ elements (nodes). This layer represents both what the RBM might be given as training input, and what we want it to be able to reconstruct. This might for example be the pixels of an image, the spin values of the Ising model, or coefficients representing speech.

\item The function $\bm{h}$ represents the hidden, or latent, layer. A vector of $N$ elements (nodes). Also called "feature detectors".
\end{enumerate}

\noindent
% !split
\subsection{Goal of hidden layer}

The goal of the hidden layer is to increase the model's expressive
power. We encode complex interactions between visible variables by
introducing additional, hidden variables that interact with visible
degrees of freedom in a simple manner, yet still reproduce the complex
correlations between visible degrees in the data once marginalized
over (integrated out).

% !split
\subsection{The parameters}

The network parameters, to be optimized/learned:
\begin{enumerate}
\item $\bm{a}$ represents the visible bias, a vector of same length $M$ as $\bm{x}$.

\item $\bm{b}$ represents the hidden bias, a vector of same length $N$  as $\bm{h}$.

\item $\bm{W}$ represents the interaction weights, a matrix of size $M\times N$.
\end{enumerate}

\noindent
Note that we have specified the lengths of $bm{x}$ and $\bm{h}$. These
lengths define the number of visible and hidden units, respectively.

% !split
\subsection{Joint distribution}

The restricted Boltzmann machine is described by a Boltzmann distribution
\begin{align*}
	P_{rbm}(\bm{x},\bm{h},\bm{\Theta}) = \frac{1}{Z(\bm{\Theta})} \exp{-(E(\bm{x},\bm{h},\bm{\Theta}))},
\end{align*}
where $Z$ is the normalization constant or partition function discussed earlier and defined as 
\begin{align*}
	Z(\bm{\Theta}) = \int \int \exp{-E(\bm{x},\bm{h},\bm{\Theta})} d\bm{x} d\bm{h}.
\end{align*}
It is common to set  the temperature $T$ to one. It is omitted in the equations above. The energy is thus a dimensionless function.

% !split

% !split
\subsection{Neural Quantum States}

The wavefunction should be a probability amplitude depending on $\bm{x}$. The RBM model is given by the joint distribution of $\bm{x}$ and $\bm{h}$
\begin{align*}
	F_{rbm}(\bm{x},\mathbf{h}) = \frac{1}{Z} e^{-\frac{1}{T_0}E(\bm{x},\mathbf{h})}
\end{align*}
To find the marginal distribution of $\bm{x}$ we set:

\begin{align*}
	F_{rbm}(\mathbf{x}) &= \sum_\mathbf{h} F_{rbm}(\mathbf{x}, \mathbf{h}) \\
				&= \frac{1}{Z}\sum_\mathbf{h} e^{-E(\mathbf{x}, \mathbf{h})}
\end{align*}

% !split
\subsection{Model for the trial wave function}
Now this is what we use to represent the wave function, calling it a neural-network quantum state (NQS)

\begin{align*}
	\Psi (\mathbf{X}) &= F_{rbm}(\mathbf{x}) \\
	&= \frac{1}{Z}\sum_{\bm{h}} e^{-E(\mathbf{x}, \mathbf{h})} \\
	&= \frac{1}{Z} \sum_{\{h_j\}} e^{-\sum_i^M \frac{(x_i - a_i)^2}{2\sigma^2} + \sum_j^N b_j h_j + \sum_{i,j}^{M,N} \frac{x_i w_{ij} h_j}{\sigma^2}} \\
	&= \frac{1}{Z} e^{-\sum_i^M \frac{(x_i - a_i)^2}{2\sigma^2}} \prod_j^N (1 + e^{b_j + \sum_i^M \frac{x_i w_{ij}}{\sigma^2}}) \\
\end{align*}

% !split
\subsection{Allowing for complex valued functions}

The above wavefunction is the most general one because it allows for
complex valued wavefunctions. However it fundamentally changes the
probabilistic foundation of the RBM, because what is usually a
probability in the RBM framework is now a an amplitude. This means
that a lot of the theoretical framework usually used to interpret the
model, i.e.~graphical models, conditional probabilities, and Markov
random fields, breaks down.

% !split
\subsection{Squared wave function}
If we assume the wavefunction to be
positive definite, however, we can use the RBM to represent the squared
wavefunction, and thereby a probability. This also makes it possible
to sample from the model using Gibbs sampling, because we can obtain
the conditional probabilities.

\begin{align*}
	|\Psi (\mathbf{X})|^2 &= F_{rbm}(\mathbf{X}) \\
	\Rightarrow \Psi (\mathbf{X}) &= \sqrt{F_{rbm}(\mathbf{X})} \\
	&= \frac{1}{\sqrt{Z}}\sqrt{\sum_{\{h_j\}} e^{-E(\mathbf{X}, \mathbf{h})}} \\
	&= \frac{1}{\sqrt{Z}} \sqrt{\sum_{\{h_j\}} e^{-\sum_i^M \frac{(X_i - a_i)^2}{2\sigma^2} + \sum_j^N b_j h_j + \sum_{i,j}^{M,N} \frac{X_i w_{ij} h_j}{\sigma^2}} }\\
	&= \frac{1}{\sqrt{Z}} e^{-\sum_i^M \frac{(X_i - a_i)^2}{4\sigma^2}} \sqrt{\sum_{\{h_j\}} \prod_j^N e^{b_j h_j + \sum_i^M \frac{X_i w_{ij} h_j}{\sigma^2}}} \\
	&= \frac{1}{\sqrt{Z}} e^{-\sum_i^M \frac{(X_i - a_i)^2}{4\sigma^2}} \sqrt{\prod_j^N \sum_{h_j}  e^{b_j h_j + \sum_i^M \frac{X_i w_{ij} h_j}{\sigma^2}}} \\
	&= \frac{1}{\sqrt{Z}} e^{-\sum_i^M \frac{(X_i - a_i)^2}{4\sigma^2}} \prod_j^N \sqrt{e^0 + e^{b_j + \sum_i^M \frac{X_i w_{ij}}{\sigma^2}}} \\
	&= \frac{1}{\sqrt{Z}} e^{-\sum_i^M \frac{(X_i - a_i)^2}{4\sigma^2}} \prod_j^N \sqrt{1 + e^{b_j + \sum_i^M \frac{X_i w_{ij}}{\sigma^2}}} \\
\end{align*}

% !split
\subsection{Cost function}

This is where we deviate from what is common in machine
learning. Rather than defining a cost function based on some dataset,
our cost function is the energy of the quantum mechanical system. From
the variational principle we know that minimizing this energy should
lead to the ground state wavefunction. As stated previously the local
energy is given by

\begin{align*}
	E_L = \frac{1}{\Psi} \hat{\mathbf{H}} \Psi.
\end{align*}

% !split
\subsection{And the gradient}

\begin{align*}
	G_i = \frac{\partial \langle E_L \rangle}{\partial \alpha_i}
	= 2(\langle E_L \frac{1}{\Psi}\frac{\partial \Psi}{\partial \alpha_i} \rangle - \langle E_L \rangle \langle \frac{1}{\Psi}\frac{\partial \Psi}{\partial \alpha_i} \rangle ),
\end{align*}
where $\alpha_i = a_1,...,a_M,b_1,...,b_N,w_{11},...,w_{MN}$.

% !split
\subsection{Additional details}

We use that $\frac{1}{\Psi}\frac{\partial \Psi}{\partial \alpha_i} 
	= \frac{\partial \ln{\Psi}}{\partial \alpha_i}$,
and find

\begin{align*}
	\ln{\Psi({\mathbf{X}})} &= -\ln{Z} - \sum_m^M \frac{(X_m - a_m)^2}{2\sigma^2}
	+ \sum_n^N \ln({1 + e^{b_n + \sum_i^M \frac{X_i w_{in}}{\sigma^2}})}.
\end{align*}

% !split
\subsection{Final equation}
This gives

\begin{align*}
	\frac{\partial }{\partial a_m} \ln\Psi
	&= 	\frac{1}{\sigma^2} (X_m - a_m) \\
	\frac{\partial }{\partial b_n} \ln\Psi
	&=
	\frac{1}{e^{-b_n-\frac{1}{\sigma^2}\sum_i^M X_i w_{in}} + 1} \\
	\frac{\partial }{\partial w_{mn}} \ln\Psi
	&= \frac{X_m}{\sigma^2(e^{-b_n-\frac{1}{\sigma^2}\sum_i^M X_i w_{in}} + 1)}.
\end{align*}

% !split
\subsection{Code example}

\textbf{This part is best seen using the jupyter-notebook}


































































































































































































































































\bpycod
# 2-electron VMC code for 2dim quantum dot with importance sampling
# Using gaussian rng for new positions and Metropolis- Hastings 
# Added restricted boltzmann machine method for dealing with the wavefunction
# RBM code based heavily off of:
# https://github.com/CompPhysics/ComputationalPhysics2/tree/gh-pages/doc/Programs/BoltzmannMachines/MLcpp/src/CppCode/ob
from math import exp, sqrt
from random import random, seed, normalvariate
import numpy as np
import matplotlib.pyplot as plt
from mpl_toolkits.mplot3d import Axes3D
from matplotlib import cm
from matplotlib.ticker import LinearLocator, FormatStrFormatter
import sys



# Trial wave function for the 2-electron quantum dot in two dims
def WaveFunction(r,a,b,w):
    sigma=1.0
    sig2 = sigma**2
    Psi1 = 0.0
    Psi2 = 1.0
    Q = Qfac(r,b,w)
    
    for iq in range(NumberParticles):
        for ix in range(Dimension):
            Psi1 += (r[iq,ix]-a[iq,ix])**2
            
    for ih in range(NumberHidden):
        Psi2 *= (1.0 + np.exp(Q[ih]))
        
    Psi1 = np.exp(-Psi1/(2*sig2))

    return Psi1*Psi2

# Local energy  for the 2-electron quantum dot in two dims, using analytical local energy
def LocalEnergy(r,a,b,w):
    sigma=1.0
    sig2 = sigma**2
    locenergy = 0.0
    
    Q = Qfac(r,b,w)

    for iq in range(NumberParticles):
        for ix in range(Dimension):
            sum1 = 0.0
            sum2 = 0.0
            for ih in range(NumberHidden):
                sum1 += w[iq,ix,ih]/(1+np.exp(-Q[ih]))
                sum2 += w[iq,ix,ih]**2 * np.exp(Q[ih]) / (1.0 + np.exp(Q[ih]))**2
    
            dlnpsi1 = -(r[iq,ix] - a[iq,ix]) /sig2 + sum1/sig2
            dlnpsi2 = -1/sig2 + sum2/sig2**2
            locenergy += 0.5*(-dlnpsi1*dlnpsi1 - dlnpsi2 + r[iq,ix]**2)
            
    if(interaction==True):
        for iq1 in range(NumberParticles):
            for iq2 in range(iq1):
                distance = 0.0
                for ix in range(Dimension):
                    distance += (r[iq1,ix] - r[iq2,ix])**2
                    
                locenergy += 1/sqrt(distance)
                
    return locenergy

# Derivate of wave function ansatz as function of variational parameters
def DerivativeWFansatz(r,a,b,w):
    
    sigma=1.0
    sig2 = sigma**2
    
    Q = Qfac(r,b,w)
    
    WfDer = np.empty((3,),dtype=object)
    WfDer = [np.copy(a),np.copy(b),np.copy(w)]
    
    WfDer[0] = (r-a)/sig2
    WfDer[1] = 1 / (1 + np.exp(-Q))
    
    for ih in range(NumberHidden):
        WfDer[2][:,:,ih] = w[:,:,ih] / (sig2*(1+np.exp(-Q[ih])))
            
    return  WfDer

# Setting up the quantum force for the two-electron quantum dot, recall that it is a vector
def QuantumForce(r,a,b,w):

    sigma=1.0
    sig2 = sigma**2
    
    qforce = np.zeros((NumberParticles,Dimension), np.double)
    sum1 = np.zeros((NumberParticles,Dimension), np.double)
    
    Q = Qfac(r,b,w)
    
    for ih in range(NumberHidden):
        sum1 += w[:,:,ih]/(1+np.exp(-Q[ih]))
    
    qforce = 2*(-(r-a)/sig2 + sum1/sig2)
    
    return qforce
    
def Qfac(r,b,w):
    Q = np.zeros((NumberHidden), np.double)
    temp = np.zeros((NumberHidden), np.double)
    
    for ih in range(NumberHidden):
        temp[ih] = (r*w[:,:,ih]).sum()
        
    Q = b + temp
    
    return Q
    
# Computing the derivative of the energy and the energy 
def EnergyMinimization(a,b,w):

    NumberMCcycles= 10000
    # Parameters in the Fokker-Planck simulation of the quantum force
    D = 0.5
    TimeStep = 0.05
    # positions
    PositionOld = np.zeros((NumberParticles,Dimension), np.double)
    PositionNew = np.zeros((NumberParticles,Dimension), np.double)
    # Quantum force
    QuantumForceOld = np.zeros((NumberParticles,Dimension), np.double)
    QuantumForceNew = np.zeros((NumberParticles,Dimension), np.double)

    # seed for rng generator 
    seed()
    energy = 0.0
    DeltaE = 0.0

    EnergyDer = np.empty((3,),dtype=object)
    DeltaPsi = np.empty((3,),dtype=object)
    DerivativePsiE = np.empty((3,),dtype=object)
    EnergyDer = [np.copy(a),np.copy(b),np.copy(w)]
    DeltaPsi = [np.copy(a),np.copy(b),np.copy(w)]
    DerivativePsiE = [np.copy(a),np.copy(b),np.copy(w)]
    for i in range(3): EnergyDer[i].fill(0.0)
    for i in range(3): DeltaPsi[i].fill(0.0)
    for i in range(3): DerivativePsiE[i].fill(0.0)

    
    #Initial position
    for i in range(NumberParticles):
        for j in range(Dimension):
            PositionOld[i,j] = normalvariate(0.0,1.0)*sqrt(TimeStep)
    wfold = WaveFunction(PositionOld,a,b,w)
    QuantumForceOld = QuantumForce(PositionOld,a,b,w)

    #Loop over MC MCcycles
    for MCcycle in range(NumberMCcycles):
        #Trial position moving one particle at the time
        for i in range(NumberParticles):
            for j in range(Dimension):
                PositionNew[i,j] = PositionOld[i,j]+normalvariate(0.0,1.0)*sqrt(TimeStep)+\
                                       QuantumForceOld[i,j]*TimeStep*D
            wfnew = WaveFunction(PositionNew,a,b,w)
            QuantumForceNew = QuantumForce(PositionNew,a,b,w)
            
            GreensFunction = 0.0
            for j in range(Dimension):
                GreensFunction += 0.5*(QuantumForceOld[i,j]+QuantumForceNew[i,j])*\
                                      (D*TimeStep*0.5*(QuantumForceOld[i,j]-QuantumForceNew[i,j])-\
                                      PositionNew[i,j]+PositionOld[i,j])
      
            GreensFunction = exp(GreensFunction)
            ProbabilityRatio = GreensFunction*wfnew**2/wfold**2
            #Metropolis-Hastings test to see whether we accept the move
            if random() <= ProbabilityRatio:
                for j in range(Dimension):
                    PositionOld[i,j] = PositionNew[i,j]
                    QuantumForceOld[i,j] = QuantumForceNew[i,j]
                wfold = wfnew
        #print("wf new:        ", wfnew)
        #print("force on 1 new:", QuantumForceNew[0,:])
        #print("pos of 1 new:  ", PositionNew[0,:])
        #print("force on 2 new:", QuantumForceNew[1,:])
        #print("pos of 2 new:  ", PositionNew[1,:])
        DeltaE = LocalEnergy(PositionOld,a,b,w)
        DerPsi = DerivativeWFansatz(PositionOld,a,b,w)
        
        DeltaPsi[0] += DerPsi[0]
        DeltaPsi[1] += DerPsi[1]
        DeltaPsi[2] += DerPsi[2]
        
        energy += DeltaE

        DerivativePsiE[0] += DerPsi[0]*DeltaE
        DerivativePsiE[1] += DerPsi[1]*DeltaE
        DerivativePsiE[2] += DerPsi[2]*DeltaE
            
    # We calculate mean values
    energy /= NumberMCcycles
    DerivativePsiE[0] /= NumberMCcycles
    DerivativePsiE[1] /= NumberMCcycles
    DerivativePsiE[2] /= NumberMCcycles
    DeltaPsi[0] /= NumberMCcycles
    DeltaPsi[1] /= NumberMCcycles
    DeltaPsi[2] /= NumberMCcycles
    EnergyDer[0]  = 2*(DerivativePsiE[0]-DeltaPsi[0]*energy)
    EnergyDer[1]  = 2*(DerivativePsiE[1]-DeltaPsi[1]*energy)
    EnergyDer[2]  = 2*(DerivativePsiE[2]-DeltaPsi[2]*energy)
    return energy, EnergyDer


#Here starts the main program with variable declarations
NumberParticles = 2
Dimension = 2
NumberHidden = 2

interaction=False

# guess for parameters
a=np.random.normal(loc=0.0, scale=0.001, size=(NumberParticles,Dimension))
b=np.random.normal(loc=0.0, scale=0.001, size=(NumberHidden))
w=np.random.normal(loc=0.0, scale=0.001, size=(NumberParticles,Dimension,NumberHidden))
# Set up iteration using stochastic gradient method
Energy = 0
EDerivative = np.empty((3,),dtype=object)
EDerivative = [np.copy(a),np.copy(b),np.copy(w)]
# Learning rate eta, max iterations, need to change to adaptive learning rate
eta = 0.001
MaxIterations = 50
iter = 0
np.seterr(invalid='raise')
Energies = np.zeros(MaxIterations)
EnergyDerivatives1 = np.zeros(MaxIterations)
EnergyDerivatives2 = np.zeros(MaxIterations)

while iter < MaxIterations:
    Energy, EDerivative = EnergyMinimization(a,b,w)
    agradient = EDerivative[0]
    bgradient = EDerivative[1]
    wgradient = EDerivative[2]
    a -= eta*agradient
    b -= eta*bgradient 
    w -= eta*wgradient 
    Energies[iter] = Energy
    print("Energy:",Energy)
    #EnergyDerivatives1[iter] = EDerivative[0] 
    #EnergyDerivatives2[iter] = EDerivative[1]
    #EnergyDerivatives3[iter] = EDerivative[2] 


    iter += 1

#nice printout with Pandas
import pandas as pd
from pandas import DataFrame
pd.set_option('max_columns', 6)
data ={'Energy':Energies}#,'A Derivative':EnergyDerivatives1,'B Derivative':EnergyDerivatives2,'Weights Derivative':EnergyDerivatives3}

frame = pd.DataFrame(data)
print(frame)

\epycod



% ------------------- end of main content ---------------

% #ifdef PREAMBLE
\end{document}
% #endif

